% Fragmento público generado desde propietarios íntegros.
% Residencia: 52.8.3.
% Fuente: ../incoming/radion_monografia_20260827/manuscrito/sections/09_sector_90120.tex:115-175
\subsubsection{Funcional hexada--octada de Barbero--Immirzi}

\paragraph{Series orientadas y salto de incidencia}

Definimos
\begin{equation}
S_{90}(q)=\frac{q^{90}}{1-q^{270}},
\qquad
S_{120}(q)=\frac{q^{120}}{1-q^{360}},
\label{eq:barbero-series}
\end{equation}
y
\[
\Delta S_m=S_m(q_-)-S_m(q_+).
\]
El funcional de inclusión marcada es
\begin{equation}
K_{6\to8}=12\Delta S_{90}-\Delta S_{120}.
\label{eq:barbero-núcleo}
\end{equation}

\paragraph{Comprobación diofántica del peso dodecafásico \(12:-1\)}

La cadena fundamental dodecafásica y su lector de retorno, construidos en
\eqref{eq:l9s102:barbero-cadena-dodecafasica}--%
\eqref{eq:l9s102:barbero-peso-generado}, producen doce términos de tipo
\(S_{90}\) y un término de frontera de tipo \(-S_{120}\). El par orientado
queda determinado antes de calcular el coeficiente de polo.

Consideremos \(aS_{90}-bS_{120}\). La comprobación del coeficiente de
\((1-q)^{-1}\) toma la forma
\begin{equation}
\frac{a}{270}-\frac{b}{360}=\frac1{24}.
\label{eq:pole-condition}
\end{equation}
Multiplicar por \(1080\) da
\[
4a-3b=45.
\]
La única costura de la vuelta fundamental tiene multiplicidad absoluta uno.
Su verificación diofántica corresponde al contratérmino entero mínimo
\(|b|=1\). Para \(b=1\), \(a=12\); para \(b=-1\), \(a=21/2\) no es entero.
Por tanto,
\begin{equation}
\boxed{a:b=12:1,\qquad aS_{90}-bS_{120}=12S_{90}-S_{120}.}
\label{eq:12minus1}
\end{equation}

\begin{proposition}[Unicidad de la comprobación diofántica]
El par generado \((12,1)\) es la única solución entera de
\eqref{eq:pole-condition} con canal principal positivo y contratérmino
mínimo \(|b|=1\).
\end{proposition}

La realización hexada--octada conserva los canales \(90/120\)
previamente construidos; la cadena dodecafásica aporta las multiplicidades
y la orientación de frontera; el lector de retorno asigna las series a los
generadores. Estas operaciones conjuntas producen \(12S_{90}-S_{120}\).
La ecuación \eqref{eq:pole-condition}, la positividad y la minimalidad
conservan su función de comprobación aritmética posterior.

La cadena Barbero es posterior:
\[
(A,C^*)\to(x,y)\to q_\pm
\to K_{6\to8}\to\gamma_{\mathrm{BI}},
\]
con evaluación
\[
\gamma_{\mathrm{BI}}=0.274007672694459\ldots.
\]
Barbero--Immirzi no genera retroactivamente \(C^*\), \(\Delta_4\) ni
\(\epsR\).

